% Compile with: pdflatex weekly_progress3.tex (run twice)

\documentclass[10pt,aspectratio=169]{beamer}

\usepackage[utf8]{inputenc}
\usepackage[T1]{fontenc}
\usepackage{lmodern}
\usepackage{amsmath, amssymb}
\usepackage{tikz}
\usepackage{booktabs}
\usepackage[table]{xcolor}
\usepackage{array}
\usepackage{graphicx}

\usetikzlibrary{positioning, shapes.geometric, arrows.meta, backgrounds, fit, calc}

\usetheme{Madrid}
\usecolortheme{seahorse}
\setbeamertemplate{navigation symbols}{}
\setbeamertemplate{caption}[numbered]
\setbeamerfont{frametitle}{size=\large}
\setbeamerfont{title}{size=\Large}

\definecolor{paperA}{RGB}{200,80,80}
\definecolor{paperB}{RGB}{60,100,180}
\definecolor{paperC}{RGB}{130,60,180}
\definecolor{paperD}{RGB}{50,140,80}
\definecolor{boxbg}{RGB}{240,240,250}
\definecolor{frozenbg}{RGB}{220,220,220}
\definecolor{trainedbg}{RGB}{255,235,205}
\colorlet{accent}{paperB}

% robust image: show file if present, else a labelled placeholder box
\newcommand{\img}[2]{\IfFileExists{#1}{\includegraphics[#2]{#1}}%
  {\fbox{\parbox[c][3cm][c]{0.6\textwidth}{\centering\ttfamily #1\\(place image here)}}}}

\title[av\_ib Weekly Progress]{Audio-Visual Information Bottleneck\\
       \normalfont\normalsize Weekly Progress --- Where the Model Looks,\\
       and What the Bottleneck Does to Attention}
\subtitle{}
\author{}
\institute{}
\date{\today}

\begin{document}

\frame{\titlepage}

% =====================================================================
\section{This week}

\begin{frame}{What I did this week}

\begin{block}{Measured where the model puts its attention (Axis 2)}
  Gradient $\times$ input saliency across 10 AVHBench clips. Video dominates the
  total attention budget ($\sim$90\,\%), but only because it has $\sim$17$\times$
  more tokens. \textbf{Per token, audio and text receive more attention than
  video.} The model is not video-biased by preference --- it is outvoted by token
  count.
\end{block}

\begin{block}{Showed the bottleneck focuses attention on the sound source}
  With the video bottleneck on, gradient saliency concentrates on the region that
  produces the sound, rather than scattering over the background. Localization
  improves: spatial entropy $0.78\!\to\!0.69$, top-5\% mass $0.53\!\to\!0.74$.
\end{block}

\begin{block}{\ldots but the model is still partly distracted}
  The bottleneck sharpens the focus, it does not make it clean: even with the IB,
  a substantial share of saliency still lands on background regions. The filter
  helps but does not fully solve grounding.
\end{block}

\end{frame}

% =====================================================================
\section{Attention budget}

\begin{frame}{Where does the model put its attention?}

\begin{columns}[T]
\begin{column}{0.55\textwidth}
\centering
\img{../slides/axis2_summary.png}{width=\textwidth}
\end{column}

\begin{column}{0.42\textwidth}

\begin{block}{Per-modality attention (10 clips)}
{\footnotesize
\renewcommand{\arraystretch}{1.15}
\begin{tabular}{l r r r}
\toprule
\textbf{Modality} & \textbf{Tokens} & \textbf{Total} & \textbf{Per tok.} \\
\midrule
Video  & $\sim$3\,500 & 89.7\,\% & 1.0$\times$ \\
Audio  &   $\sim$200  &  6.2\,\% & \textbf{1.16}$\times$ \\
Text   &    $\sim$30  &  4.1\,\% & \textbf{5.1}$\times$ \\
\bottomrule
\end{tabular}}
\end{block}

\begin{beamercolorbox}[sep=5pt,rounded=true,colsep*=0pt]{block body}
{\footnotesize
\textbf{Takeaway.} Video has 17$\times$ more tokens than audio, so its bulk
saliency dominates. But \emph{per token}, audio and text receive more attention
than video. Compressing video with the bottleneck brings the effective token
budget closer to audio --- without any explicit reweighting.}
\end{beamercolorbox}

\end{column}
\end{columns}

\end{frame}

% =====================================================================
\begin{frame}{How we measure attention (clean, causal)}

\begin{columns}[T]
\begin{column}{0.55\textwidth}
We hook the first decoder layer and backpropagate through the log-probability of
the model's own answer token:
\[ s_i = \mathrm{ReLU}\!\Big(\nabla_{x_i}\,\log p(\hat y)\;\cdot\;x_i\Big) \]
\begin{itemize}
  \item Causal --- measures what actually drives the output.
  \item Immune to attention sinks (unlike raw decoder attention).
  \item Same checkpoint, bottleneck on vs.\ off --- a fair comparison.
\end{itemize}
\end{column}

\begin{column}{0.43\textwidth}
\begin{beamercolorbox}[sep=5pt,rounded=true,colsep*=0pt]{block body}
{\footnotesize
\textcolor{paperB}{\textbf{Why it matters.}} Raw decoder attention is dominated
by sinks and looks like noise. Gradient $\times$ input saliency yields clean,
object-aligned maps --- and lets us compare the model with and without the
bottleneck on the same trained weights.}
\end{beamercolorbox}
\end{column}
\end{columns}

\end{frame}

% =====================================================================
\section{Bottleneck localization}

\begin{frame}{The bottleneck focuses on the sound source --- but is still distracted}

\begin{center}
\img{grad_map.png}{width=0.92\textwidth}
\end{center}

\vspace{2pt}
\begin{beamercolorbox}[sep=5pt,rounded=true,colsep*=0pt]{block body}
{\footnotesize
\textcolor{paperB}{\textbf{Takeaway.}} \emph{Original / Without IB / With IB} for
two clips of a person spraying water (the hose is the sound source). Without the
bottleneck (middle), saliency scatters across the whole frame --- foliage,
background, edges. With the bottleneck (right), the strongest saliency moves onto
the arm and the spray nozzle where the sound originates. \textbf{But the focus is
not clean}: hot spots still remain on the background. Localization improves
quantitatively --- spatial entropy $0.78\!\to\!0.69$, top-5\% mass
$0.53\!\to\!0.74$ --- yet the model is only \emph{partly} grounded on the
sound-producing region.}
\end{beamercolorbox}

\end{frame}

% =====================================================================
\begin{frame}{Localization metrics --- the bottleneck sharpens, not solves}

\begin{columns}[T]
\begin{column}{0.50\textwidth}

\begin{block}{Saliency-map statistics (6 clips)}
{\footnotesize
\renewcommand{\arraystretch}{1.2}
\begin{tabular}{l r r r}
\toprule
\textbf{Metric} & \textbf{No IB} & \textbf{With IB} & \textbf{$\Delta$} \\
\midrule
Spatial entropy & 0.78 & \textbf{0.69} & $-11$\,\% \\
Top-5\% mass    & 0.53 & \textbf{0.74} & $+40$\,\% \\
\bottomrule
\end{tabular}}
\vspace{4pt}
{\footnotesize
\textbf{Spatial entropy}: $H=-\sum_i \bar s_i \log \bar s_i$ of the normalised
saliency map (lower = more localised).\\
\textbf{Top-5\% mass}: fraction of total saliency in the top 5\% of positions
(higher = more concentrated).}
\end{block}

\end{column}

\begin{column}{0.47\textwidth}

\begin{beamercolorbox}[sep=5pt,rounded=true,colsep*=0pt]{block body}
{\footnotesize
The bottleneck moves both metrics in the right direction: attention becomes more
concentrated and lower-entropy. But top-5\% mass is still only 0.74 --- a quarter
of the saliency remains spread over the rest of the frame.}
\end{beamercolorbox}

\vspace{6pt}

\begin{block}{Why the residual distraction matters}
{\footnotesize
A clean grounding on the sound source would be the strongest evidence that the
bottleneck builds a fail-safe representation: if attention sits only on the
sound-producing region, corrupting that region should reliably collapse the
answer. The residual background attention is exactly the leak we still need to
close.}
\end{block}

\end{column}
\end{columns}

\end{frame}

% =====================================================================
\section{Audio attention}

\begin{frame}{Audio attention: which segments drive the answer}

\begin{columns}[T]
\begin{column}{0.55\textwidth}

\begin{block}{Method}
{\footnotesize
Same gradient saliency pass, restricted to audio token positions. We map each
audio token back to its time offset in the clip and overlay the saliency curve on
the waveform. Top-$k$ attended time windows are marked.}
\end{block}

\begin{block}{Observations (30 AVHBench clips)}
{\footnotesize
\begin{itemize}\setlength\itemsep{2pt}
  \item The peak attended audio segment corresponds to the onset of the relevant
        sound event in $\sim$70\,\% of clips.
  \item Audio saliency is bursty: a 1--2\,s window typically accounts for $>$50\,\%
        of total audio saliency.
  \item The bottleneck model and the vanilla model show similar audio saliency
        patterns --- confirming the bottleneck acts on video, not audio.
\end{itemize}}
\end{block}

\end{column}

\begin{column}{0.42\textwidth}

\begin{block}{What this rules out}
{\footnotesize
If audio saliency were uniform, the model would be integrating the whole waveform
regardless of content. The bursty pattern shows the model attends to
\emph{events} in the audio stream, not just the presence or absence of sound.}
\end{block}

\begin{block}{Consistency with the budget analysis}
{\footnotesize
Audio saliency is unchanged by the bottleneck, matching the per-token finding:
audio tokens were always attended carefully --- the bottleneck only removes the
numerical dominance of video bulk.}
\end{block}

\end{column}
\end{columns}

\end{frame}

% =====================================================================
\section{Next steps}

\begin{frame}{Next steps}

\begin{block}{1.\ Close the residual-distraction gap}
  The bottleneck sharpens attention but a quarter of the saliency still lands on
  the background. Test whether stronger compression (higher $\beta$) pushes
  top-5\% mass higher, and whether that trades off against accuracy.
\end{block}

\begin{block}{2.\ Tie localization to fail-safety}
  If attention concentrates on the sound source, corrupting that source should
  collapse the answer. Measure the correlation between top-5\% mass and the drop
  in answer confidence under video corruption.
\end{block}

\begin{block}{3.\ Per-clip audio--video grounding}
  Combine the audio time-saliency and the video spatial-saliency on the same
  clips to check the model attends to the \emph{same event} in both streams ---
  the signature of genuine audio-visual grounding.
\end{block}

\end{frame}

\end{document}
